\chapter{کلیات تحقیق}
\label{ch:introduction}

\section{بیان مسئله}
بخش مراقبت‌های ویژه (\lr{ICU}) یکی از داده‌محورترین و در عین حال پرخطرترین محیط‌های نظام سلامت است. در ساعات نخست بستری، داده‌های آزمایشگاهی، علائم حیاتی، سن، مسیر پذیرش و تشخیص‌های اولیه به‌تدریج تصویری از وضعیت بیمار می‌سازند. پژوهش‌های یادگیری ماشین می‌کوشند از همین داده‌های اولیه برای تخمین خطر، پیامد، و مسیر احتمالی بیمار استفاده کنند. بااین‌حال، هر مدل بالینی باید با دقت تعریف کند که دقیقاً چه چیزی را پیش‌بینی می‌کند و از چه چیزی نمی‌توان نتیجه گرفت.

این پایان‌نامه از ایدۀ «فیوشات درمانی» آغاز شد: آیا می‌توان با استفاده از توکن‌های علائم و یادگیری چندنمونه‌ای (\lr{Few-Shot Learning}) از داده‌های محدود و حساس بالینی مدلی ساخت که به پیش‌بینی امکان‌پذیری درمان کمک کند؟ در اجرای نهایی و بازنگری‌شده، این ایده به یک مسئلۀ قابل‌بازتولید و محدودتر تبدیل شد: پیش‌بینی \emph{نشانگر جانشین امکان‌پذیری ترخیص از \lr{ICU}} در پایگاه دادۀ \lr{MIMIC\text{-}IV v3.1}. این نشانگر جایگزین قضاوت پزشک، آمادگی واقعی ترخیص، یا توصیه درمان نیست. هدف، ارزیابی پژوهشی مدل‌ها روی یک برچسب گذشته‌نگر و شفاف است.

\section{تعریف عملیاتی امکان‌پذیری}
در این پژوهش، برچسب مثبت زمانی به بیمار داده می‌شود که سه شرط برقرار باشد: بیمار تا ترخیص بیمارستانی زنده مانده باشد؛ مقصد ترخیص در گروه مقصدهای مطلوب تعریف‌شده باشد؛ و طول اقامت در \lr{ICU} کمتر از \pnum{720} ساعت باشد. این تعریف از روی داده‌های ساخت‌یافتۀ \lr{MIMIC\text{-}IV} قابل استخراج است و امکان بازتولید دارد، اما یک برچسب جانشین است. بنابراین، در سراسر پایان‌نامه از اصطلاح «نشانگر جانشین امکان‌پذیری ترخیص از \lr{ICU}» استفاده می‌شود.

\section{انگیزه و اهمیت}
تصمیم‌های مرتبط با مسیر مراقبت و ترخیص در \lr{ICU} پیامدهای مستقیم برای بیمار، تخت‌های مراقبت ویژه، هزینه و برنامه‌ریزی بیمارستان دارند. اگر مدلی بتواند از داده‌های \pnum{24} ساعت نخست سیگنالی قابل‌اعتماد درباره پیامد جانشین تولید کند، می‌تواند در آینده به پژوهش‌های پشتیبان تصمیم، اولویت‌بندی پایش، یا طراحی مطالعات آینده‌نگر کمک کند. اما چنین مدلی تنها زمانی ارزشمند است که با خط پایه‌های قوی مقایسه شود، کالیبراسیون آن بررسی گردد، و محدودیت‌های برچسب و داده آشکارا بیان شود.

انگیزۀ دوم، حریم خصوصی و محدودیت دسترسی به داده‌های سلامت است. پایگاه‌هایی مانند \lr{MIMIC\text{-}IV} با وجود ارزش پژوهشی بالا، طبق مقررات سخت‌گیرانه و پس از احراز صلاحیت در دسترس قرار می‌گیرند \cite{johnson2023mimic,goldberger2000physionet}. در دنیای واقعی، داده‌های بیمارستانی معمولاً میان مراکز قابل‌انتقال نیستند. به همین دلیل، این پایان‌نامه علاوه بر مدل‌های متمرکز، یک شبیه‌سازی یادگیری فدرال (\lr{Federated Learning}) را نیز بررسی می‌کند \cite{mcmahan2017communication,li2019federated}.

\section{شکاف پژوهشی}
مطالعات زیادی نشان داده‌اند که مدل‌های یادگیری ماشین می‌توانند از داده‌های \lr{EHR} برای پیش‌بینی پیامدهای بالینی استفاده کنند \cite{rajkomar2018scalable}. بااین‌حال، سه شکاف در این حوزه باقی است. نخست، بسیاری از مطالعات، خط پایه‌های کلاسیک قوی و مدل‌های نوین را در یک چارچوب یکسان و با گزارش کالیبراسیون مقایسه نمی‌کنند. دوم، یادگیری چندنمونه‌ای در داده‌های جدولی ساخت‌یافته، برخلاف تصویر و زبان، همیشه مزیت روشن ندارد و نیازمند ارزیابی صادقانه است \cite{snell2017prototypical,vinyals2016matching}. سوم، کاربرد بازنمایی‌های توکنی علائم در کنار مدل‌های جدولی و شبیه‌سازی فدرال هنوز نیازمند بنچمارک‌های بازتولیدپذیر است.

\section{پرسش‌های پژوهش}
این پایان‌نامه به پنج پرسش پاسخ می‌دهد:
\begin{enumerate}
\item آیا ویژگی‌های آزمایشگاهی و علائم حیاتی \pnum{24} ساعت نخست \lr{ICU} می‌توانند نشانگر جانشین امکان‌پذیری ترخیص را در \lr{MIMIC\text{-}IV} پیش‌بینی کنند؟
\item مدل‌های کلاسیک و درختی جدولی در مقایسه با مدل \lr{Few-Shot ETHOS} چگونه عمل می‌کنند؟
\item آیا بازنمایی توکن‌های علائم، به‌ویژه در مدل‌های هیبریدی، سیگنال مکمل فراهم می‌کند؟
\item آیا شبیه‌سازی فدرال می‌تواند عملکردی نزدیک به خط پایۀ متمرکز هم‌خانواده حفظ کند؟
\item کدام مدل بهترین توازن میان تفکیک (\lr{AUROC/AUPRC}) و کالیبراسیون (\lr{Brier/ECE}) را دارد؟
\end{enumerate}

\section{مشارکت‌ها}
مشارکت‌های اصلی پژوهش عبارت‌اند از:
\begin{enumerate}
\item ساخت و اعتبارسنجی یک کوهورت کامل \lr{MIMIC\text{-}IV} شامل \pnum{32,399} بستری نخست بزرگسال در \lr{ICU}.
\item تعریف شفاف یک نشانگر جانشین امکان‌پذیری ترخیص از \lr{ICU} و تفکیک آن از آمادگی واقعی ترخیص یا توصیه درمان.
\item ساخت بازنمایی توکن‌های علائم برای داده‌های آزمایشگاهی و علائم حیاتی \pnum{24} ساعت نخست.
\item مقایسۀ مدل‌های رگرسیون لجستیک، جنگل تصادفی، گرادیان‌بوستینگ، خط پایۀ عصبی، \lr{Few-Shot ETHOS}، مدل هیبریدی، مدل انباشتی، و شبیه‌سازی فدرال.
\item نشان دادن اینکه در این وظیفه، مدل‌های درختی و جدولی از \lr{Few-Shot ETHOS} خام قوی‌ترند، در حالی که توکن‌های علائم در حالت هیبریدی می‌توانند به عملکرد نزدیک به مدل‌های قوی برسند.
\item گزارش معیارهای کالیبراسیون در کنار معیارهای تفکیک، و انتخاب \lr{XGBoost} به‌عنوان مدل اصلی قابل‌دفاع‌تر از نظر توازن عملکرد.
\end{enumerate}

\section{ساختار پایان‌نامه}
فصل دوم پیشینۀ علمی را مرور می‌کند. فصل سوم روش‌شناسی، داده، برچسب، ویژگی‌ها، توکن‌ها، مدل‌ها، تقسیم داده و معیارهای ارزیابی را توضیح می‌دهد. فصل چهارم نتایج کامل کوهورت، جداول و شکل‌ها را گزارش می‌کند. فصل پنجم یافته‌ها، محدودیت‌ها و پیامدها را بحث می‌کند و فصل ششم جمع‌بندی نهایی و مسیرهای آینده را ارائه می‌دهد.
